\documentclass[12pt]{article}
\usepackage[utf8]{inputenc}
\usepackage[T1]{fontenc}
\usepackage{hyperref}
\title{Southern Tiwa Language Resources for K-12 Curriculum Development}
\author{Compiled by Claude Sonnet 5 (claude-sonnet-5) for Andr\'as Kornai}
\date{}
\begin{document}
\maketitle
\section{Introduction}
Southern Tiwa (ISO 639-3 tix) is a Kiowa-Tanoan language of New Mexico (Isleta and Sandia Pueblos) and Texas (Ysleta del Sur). This search found little that is usable for a curriculum. Published classroom materials were not found, and what exists is mostly reference and archival.

\section{Primary Educational Resources}
No primer, textbook or school programme for Southern Tiwa turned up in the searches.

\section{Linguistic and Reference Materials}
\begin{itemize}
\item Wikipedia overview of Southern Tiwa: \url{https://en.wikipedia.org/wiki/Southern_Tiwa_language}
\item Omniglot page with an alphabet and pronunciation chart: \url{https://omniglot.com/writing/southerntiwa.htm}
\item Smithsonian National Anthropological Archives finding aid (MS 1027), which covers copied Isleta vocabularies and a comparative Tiwa vocabulary: \url{https://sirismm.si.edu/EADpdfs/NAA.MS1027.pdf}
\item Endangered Languages Project entry: \url{https://www.endangeredlanguages.com/node/85513}
\end{itemize}
Cited in the searched sources, but not located online: W. L. Leap, \emph{The language of Isleta, New Mexico} (1970 dissertation), and George Trager's papers on Tiwa phonology and grammar.

\section{Religious and Community Resources}
None found. native-languages.org states that many Tiwa elders hold that the language is not meant to be written, and that site deliberately omits vocabulary and texts: \url{https://www.native-languages.org/tiwa.htm}

\section{Digital and Audio Resources}
No audio, video or software resources were found.

\section{Curriculum Development Context}
Any curriculum work would need to be led by, and agreed with, the Pueblo communities, given the cultural restrictions on writing the language. The sources found here do not support building a curriculum on their own.
\end{document}